% Modernised reading — fonds Alexandre Grothendieck, shelfmark 29, batch 5
% (pages 81-100).
%
% This is an INTERPRETATION, not a transcription. It re-reads the pages in
% today's notation and vocabulary, states the hypotheses the manuscript leaves
% implicit, and drops the critical apparatus so that the mathematics reads
% straight through. Everything the brackets and underlines carried moves into a
% footnote. In any dispute of substance the transcription governs: whoever
% cites Grothendieck must cite batch-05.fr.tex.
%
% Scope: the folder was read WHOLE before any of this was written — all eleven
% batches, pages 1-216, in one pass — and one file was then written per batch.
% The one observation this file could not have made batch by batch is the
% central one: the theorem of page 89 (the extension class is the Chern class
% of the self-intersection), the programme of page 100 (the same for a general
% family of divisors), the redraft of page 102, and the parenthesis of page 200
% (the obstruction to a section of the Kummer gerbe is the class of the D_i)
% are FOUR statements of one theorem, spread over a hundred and ten pages.
%
% Departures from the page, each footnoted where it occurs:
%   - page 89 puts the class of a line bundle in H^1(X_o, mu) where page 95
%     puts the same class in H^2; the two are related by the Kummer sequence
%     and both writings are recorded;
%   - page 93 asserts that the kernel of Pic(X_m) -> Pic(X_o) is uniquely
%     n-divisible without saying why. The reason — n is invertible on X_o and
%     the kernel is filtered by cohomology of O_{X_o}-modules — is supplied;
%   - page 97 asserts mu^infty = 0 for X_o = Spec Z without saying why. The
%     reason is supplied.
%
% Pass: Opus 5 (claude-opus-5), 2026-08-29 — first pass, unchecked against the
% pages by a human.
\documentclass[11pt,a4paper]{article}
% Le préambule commun à toutes les transcriptions du fonds Grothendieck.
%
% Il tient en une idée : **l'incertitude se compose, elle ne se résout pas.**
% Une transcription qui lisse un mot douteux a détruit la seule chose pour
% laquelle on la faisait — un lecteur doit pouvoir savoir, à chaque endroit,
% ce qui était sur le feuillet et ce qui a été ajouté. Les six macros qui
% suivent sont donc l'appareil critique, et non de la décoration.
%
% Les mêmes commandes sont reconnues par `scripts/render.mjs`, qui produit la
% vue de lecture du site. Une macro ajoutée ici doit l'être aussi là-bas,
% faute de quoi le rendu échouera bruyamment — ce qui est voulu : un rendu
% silencieusement faux est pire qu'un rendu absent.

\ProvidesPackage{grothendieck}[2026/08/08 Transcriptions du fonds Grothendieck]

\usepackage{amsmath,amssymb,amsthm}
\usepackage{mathrsfs}
\usepackage{stmaryrd}
% Les diagrammes commutatifs sont partout dans ce fonds : c'est la langue
% même de la géométrie algébrique de Grothendieck, et une transcription qui
% les rendrait en prose ne transcrirait rien.
\usepackage{tikz-cd}
% Français par défaut — c'est la langue des feuillets — mais l'anglais est
% chargé aussi : la traduction et le résumé sont des documents anglais, et la
% typographie française (espaces avant les deux-points) y serait une faute.
% Ces documents basculent par \selectlanguage{english} en tête de corps.
\usepackage[english,french]{babel}
\usepackage{fontspec}
\usepackage{geometry}
\geometry{a4paper,margin=25mm,marginparwidth=28mm}
\usepackage{marginnote}
\usepackage{xcolor}
% ulem, not soul. `soul`'s \st and \ul are built on a character-by-character
% scan that XeLaTeX's font machinery breaks: the document fails with an
% undefined control sequence rather than degrading. ulem does the same two jobs
% and is Unicode-engine safe. `normalem` stops it hijacking \emph, which the
% transcriptions use for Grothendieck's own emphasis.
\usepackage[normalem]{ulem}
\usepackage{hyperref}
\hypersetup{colorlinks=true,linkcolor=black,urlcolor=black,pdfborder={0 0 0}}

% --- Métadonnées du lot -----------------------------------------------------
% Reprises telles quelles par le script de rendu, qui les affiche en tête de
% la vue de lecture. Elles fixent ce que le fichier prétend couvrir : sans
% elles, une transcription flotte sans point d'attache dans un fonds de seize
% mille feuillets.

\newcommand{\folder}[1]{\gdef\grothFolder{#1}}
\newcommand{\batch}[1]{\gdef\grothBatch{#1}}
\newcommand{\leaves}[2]{\gdef\grothFirst{#1}\gdef\grothLast{#2}}
\newcommand{\foldertitle}[1]{\gdef\grothFoldertitle{#1}}
\gdef\grothFolder{?}\gdef\grothBatch{?}\gdef\grothFirst{?}\gdef\grothLast{?}\gdef\grothFoldertitle{}

% --- Le résumé --------------------------------------------------------------
%
% Il ouvre la lecture modernisée, sous le titre « Résumé », et il est composé
% en petit. Ce n'est pas un ornement : c'est le seul passage du document qui
% ne soit pas des mathématiques, et un lecteur doit pouvoir le reconnaître
% avant de l'avoir lu — puis le sauter s'il connaît déjà le sujet, ou s'y
% arrêter s'il ne le connaît pas. Le filet qui le suit marque la reprise du
% corps.

\newenvironment{resume}
  {\small\setlength{\parskip}{0.45em}}
  {\par\medskip\hrule\medskip}

% --- Mots-clés ---------------------------------------------------------------
%
% En anglais, en clôture du résumé de la lecture modernisée : le vocabulaire
% moderne sous lequel chercher ce que les feuillets construisent. Le manifeste
% du site (`npm run manifest`) les extrait pour étiqueter la cote entière —
% cette ligne est la source unique des tags, il n'y a pas de fichier à côté.
\newcommand{\keywords}[1]{\par\smallskip\noindent
  {\footnotesize\textsc{Keywords} --- \textit{#1}}\par}

% --- L'appareil critique ----------------------------------------------------

% Le numéro de feuillet, en marge. C'est l'ancre qui permet de régler un
% désaccord sur une formule en regardant *un* feuillet plutôt qu'une chemise
% de six cents. Le site s'en sert aussi pour tourner les pages du fac-similé
% quand on fait défiler la transcription.
\newcommand{\leaf}[1]{%
  \marginnote{\footnotesize\color{gray!70}\textsf{#1}}%
  \label{leaf:#1}\ignorespaces}

% Illisible : on ne devine pas. Un mot inventé qui se lit comme les autres est
% le pire résultat possible.
\newcommand{\ill}{\textcolor{red!60!black}{[\,\ldots\,]}}

% Lecture douteuse : le mot est donné, et signalé comme incertain.
\newcommand{\uncertain}[1]{\uline{#1}}

% Ajout de l'éditeur — une lettre manquante, un mot sous-entendu.
\newcommand{\add}[1]{\textcolor{blue!50!black}{[#1]}}

% Biffé par Grothendieck lui-même. Ce qu'il a rayé fait partie du document :
% une rature dit souvent où la pensée a changé de direction.
\newcommand{\struck}[1]{\sout{#1}}

% Note du transcripteur, sur l'état matériel du feuillet.
\newcommand{\note}[1]{\par\smallskip{\small\color{gray!60!black}\textit{[#1]}}\par\smallskip}

% Note marginale de Grothendieck, distincte du corps du texte.
\newcommand{\marginal}[1]{\par\smallskip\noindent\hspace{1em}%
  {\small\color{gray!60!black}#1}\par\smallskip}

% --- Titre ------------------------------------------------------------------

\AtBeginDocument{%
  \begin{center}
    {\large\bfseries Fonds Alexandre Grothendieck --- cote n\textsuperscript{o}~\grothFolder}\\[2pt]
    {\itshape\grothFoldertitle}\\[4pt]
    {\small feuillets \grothFirst--\grothLast\ (lot \grothBatch)}\\[2pt]
    {\footnotesize\color{gray!60!black}Transcription établie d'après le fac-similé numérisé
     par l'Université de Montpellier.\\
     Les lectures incertaines sont soulignées, les illisibles notées [\,\ldots\,],
     les ajouts entre crochets.}
  \end{center}
  \vspace{1em}}

\endinput


\folder{29}
\batch{5}
\pages{81}{100}
\dating{1959-1969}
\watermark{Édition de démonstration\\interprétation personnelle de l'œuvre}
\foldertitle{Groupe fondamental [Autour de SGA 4 et SGA 7] — lecture modernisée}

\begin{document}

\section*{Résumé}

\begin{resume}
Le cœur de ce lot est un calcul, et il faut savoir d'avance ce qu'il calcule.

Imaginez une famille de courbes qui dégénère : au-dessus d'un petit disque, une
courbe lisse pour chaque point sauf le centre, où elle se casse. Ce qui
intéresse, c'est ce que la famille fait quand on fait le tour du centre sans
jamais y passer : les points de la courbe reviennent permutés, et cette
permutation --- la \emph{monodromie} --- est le seul invariant qui distingue
une dégénérescence d'une autre. La question de ces pages est : de quoi
dépend-elle exactement ?

La réponse est d'une netteté remarquable, et elle tient en une phrase : elle
dépend de la manière dont la fibre centrale se coupe elle-même. Une courbe
plongée dans une surface a un nombre d'auto-intersection, entier, qui mesure
comment elle se déplace dans la surface qui la contient ; c'est le même entier
qui gouverne la monodromie. Le lemme 2 de la page 89 le dit exactement : la
classe de l'extension de groupes fondamentaux est \emph{la classe de Chern de
l'auto-intersection de la fibre spéciale}. Et l'exemple final le rend
palpable : si la fibre centrale est une droite projective d'auto-intersection
$-d$, alors la monodromie modérée est le groupe des racines $n$-ièmes de
l'unité, $n$ étant la partie de $d$ première à la caractéristique.

Ce théorème n'est pas isolé dans le carton. La page 100 le pose en programme
pour une famille quelconque de diviseurs, la page 102 le rédige en corollaire,
et la page 200 --- deux cents pages plus loin, dans une chemise qui n'a plus
rien à voir --- le retrouve sous une autre forme : l'obstruction à trancher une
racine $n$-ième d'un diviseur est la classe de ce diviseur. C'est la même
identité écrite quatre fois. Il n'y a pas dans ce carton de théorème plus
souvent réénoncé, et il faut voir dans cette insistance ce qu'elle est : la
reconnaissance progressive que le nombre d'auto-intersection et la monodromie
sont deux noms d'une seule quantité.

Trois feuillets encadrent le calcul et méritent d'être lus. Les notes de
lecture de Grothendieck sur le texte de Murre --- en anglais, chapitre par
chapitre --- contiennent une ligne qui, à cette distance, se lit comme un
programme : « Appl.\ descent for tamely ramified coverings ? At least one would
like to have a \emph{champ} on $\mathrm{Et}(S)$ ! ». C'est mot pour mot ce que
les pages 190 et suivantes construiront. Et deux feuillets de son propre
tapuscrit portent, corrigée à la main, la phrase qui baptise l'\emph{extension
de Raynaud}.

Les noms modernes à chercher sont \emph{suite exacte de spécialisation},
\emph{dégénérescence semi-stable}, \emph{fibré normal}, \emph{théorie de
Kummer} et --- pour ce que les notes de lecture appellent de leurs vœux ---
\emph{champ}.

\keywords{specialization exact sequence, log smooth degeneration, self-intersection, Chern class of a divisor, Raynaud extension}
\end{resume}

\pagerange{81}{85}
\section*{Notes de lecture, et l'appel d'un champ (pages 81, 83 et 85)}

Trois feuillets, en anglais, notent chapitre par chapitre ce que Grothendieck
veut voir changer dans le texte de Murre --- celui que la lettre du 30 avril
1969, page 78, accompagnait. Ce sont des notes de travail, faites de renvois et
de phrases coupées, et il n'y a pas lieu de les paraphraser toutes. Six items,
qui disent chacun quelque chose.

\begin{itemize}
  \item \emph{Se restreindre aux morphismes adiques.} L'ambiguïté sur ce qu'est
    un point d'un schéma formel est jugée « quite general » --- et le chapitre
    IV se voit prescrire d'\emph{introduire les points géométriques} de
    $\mathcal{X}$, plutôt que de raisonner sur
    $\operatorname{Spec}(\mathcal{O}_{\mathcal{X},s})$.
  \item \emph{Le chapitre III ne sert qu'à un énoncé.} « whole ch.\ III seems
    used only to get 3.5.8 ». Et il faudrait y ajouter deux choses : une
    interprétation topologique de $\pi_{1}^{D}$, et une section sur le cas où
    les diviseurs \emph{ont des équations} --- « This should make the situation
    much clearer ! ». C'est le cas particulier a) du lot 1, et c'est en effet
    celui où la donnée kummérienne se décrit sans détour.
  \item \emph{Le champ.} « Appl.\ descent for tamely ramified coverings ? At
    least one would like to have a \emph{champ} on
    $\mathrm{Et}(S)$ ! »\footnote{Le mot est en français dans ces notes
    anglaises, comme il l'est dans le tapuscrit de 1967 (page 22 : « qui est
    même un champ »). C'est exactement l'objet que les pages 190 à 200 de ce
    même carton définiront sous le nom de \emph{domaine de ramification} : un
    champ sur $\mathrm{Et}(S)$, muni d'une réalisation géométrique. Le carton
    ne date ni ces notes ni ces pages, et cette lecture n'en tire aucune
    chronologie ; le vœu et sa réalisation sont l'un et l'autre dans le même
    carton, à cent pages de distance.}
  \item \emph{La normalité est un artifice.} « Normality-assumption unnatural
    (formal normality, as formal schemes\dots) ». C'est le même reproche qu'il
    faisait à Murre en 1967 (page 8 : « I am convinced that these assumptions
    are anyhow artificial »), et que Murre lui retourne dans ses commentaires
    (page 79 : la normalité n'est peut-être pas stable par localisation étale).
  \item \emph{Où est le lemme d'Abhyankar ?} La question est posée deux fois,
    page 81 et en marge de la page 83 : « yoga Abhyankar lemma to be used ».
    C'est la reconnaissance que le résultat qui, dans le cas des croisements
    normaux, transforme la structure en propriété, n'a pas trouvé sa place dans
    la rédaction formelle.
  \item \emph{Le dernier item, page 85, est le plus concret}, et c'est un
    diagnostic de rédaction : ce que le chapitre X démontre longuement revient
    à montrer qu'un revêtement de $S$ est modéré relativement à $D_{1}, D_{2}$
    si et seulement si ses restrictions à $S_{1}$, $S_{2}$ le sont ---
    « \emph{should be trivial} with the conditions considered ! ».
\end{itemize}

\subsection*{Deux feuillets de son propre tapuscrit : l'extension de Raynaud}

Les pages 82 et 84 sont deux feuillets d'un tapuscrit français à lui, dans le
désordre --- ses pages 74 et 69 ---, portant des corrections de sa main. Ils
n'ont pas de rapport direct avec ce qui précède, mais l'un d'eux baptise un
objet, et le baptême est fait à la main sur la page.

La page 84 démontre que, dans le cas favorable où le schéma formel abélien
$\hat{B}$ est algébrisable et où $T_{o}$ est isotrivial, on trouve un schéma en
groupes $A^{\sharp}$ sur $S$, extension canonique
\[
  0 \longrightarrow T \longrightarrow A^{\sharp} \longrightarrow B
  \longrightarrow 0
\]
du schéma abélien $B$ par le tore $T$ qui relève $T_{o}$, caractérisée à
isomorphisme unique près par le fait de donner un système cohérent
d'isomorphismes d'extensions $A^{\sharp}_{n} \simeq A_{n}$, c'est-à-dire un
isomorphisme de schémas formels en groupes
$(A^{\sharp})^{\wedge} \simeq \hat{A}$ respectant les structures d'extension.
Elle dépend fonctoriellement de $A$ et sa formation commute au changement de
base.

La phrase centrale a été refaite à la main, et le carton conserve les deux
états : le tapuscrit annonçait « une extension canonique du schéma abélien $B$
\dots\ que nous appellerons l'\emph{extension de Raynaud} associée à $A$ », et
l'insertion manuscrite substitue « un schéma en groupes $A^{\sharp}$ sur $S$
appelé le \emph{groupe de Raynaud} associé au schéma en groupes $A$ sur
$S_{o}$ ». Le nom se déplace de l'extension à l'objet, sur la
page.\footnote{La page 82 --- sa page 74, donc postérieure dans le texte ---
porte le lemme qui sert : les extensions de $\underline{\mathbb{Z}}$ par
$G^{o\sharp}$ correspondent aux torseurs sous
$H = \underline{\mathrm{Hom}}(\underline{\mathbb{Z}}, G^{o\sharp})$, et,
$H$ étant fini localement libre sur $S$, la catégorie des torseurs sous $H$
équivaut par complétion formelle à celle des systèmes cohérents de torseurs
sous les $H_{n}$. La transcription signale que le caractère lu
$\underline{\mathbb{Z}}$ est douteux ; la lecture qui suit n'en dépend pas.}

\pagerange{86}{86}
\section*{Images inverses et voisinages infinitésimaux (page 86)}

Un feuillet d'une écriture très fine, sur un papier que traverse le tapuscrit
du verso ; la plus grande partie n'en est pas lisible et cette lecture n'y
supplée pas. Ce qui s'en dégage est une question, non un énoncé : l'image
inverse d'un module quasi-cohérent est quasi-cohérente, celle d'un module plat
est plate, et il faut vérifier que cela respecte la notion voulue ; puis, si
$X' \to X \times_{S} S'$ est non ramifié, il en va de même du morphisme induit
sur les espaces tangents, et l'on peut passer aux voisinages infinitésimaux de
tous ordres.

\begin{tikzcd}
X \times_{S} S' \arrow[r] \arrow[d] & X \arrow[d] \\
S' \arrow[r] & S
\end{tikzcd}

Le cube au crayon dont ce carré n'est qu'une face n'est pas reproduit : la
transcription en donne la raison, qui est la bonne --- le dessiner reviendrait
à affirmer ce qu'on n'a pas lu.

\pagerange{87}{89}
\section*{Le $\pi_{1}$ modéré d'un schéma formel le long de sa fibre spéciale (pages 87 et 89)}

Voici le calcul. Grothendieck pagine cette suite 1 à 6.

\subsection*{La situation}

Soit $\mathcal{X}$ un schéma formel régulier connexe, $X_{o} = \mathcal{X}_{o}$
un diviseur régulier --- la fibre spéciale ---, et
\[
  D \;=\; \bigcup_{i=0}^{n} D_{i}
\]
un diviseur à croisements normaux, les $D_{i}$ réguliers irréductibles, avec
$D_{o} = X_{o}$. On pose $D' = \sum_{i \neq 0} D_{i}$, on note $D'_{o}$ le
diviseur induit sur $X_{o}$, et $\mathcal{U} = \mathcal{X} - D_{o}$.

Il faut avoir en tête le dessin que la marge de la page porte : $D_{o} = X_{o}$
horizontal, coupé transversalement par $D_{1}$ et $D_{2}$, et le point de base
$\bar{\xi}_{o}$ sur $X_{o}$ en dehors des croisements.

\subsection*{La suite exacte}

Elle est double, et la seconde ligne dévisse le noyau de la première :
\[
  e \longrightarrow I_{o} \longrightarrow
  \pi_{1}^{t}(\mathcal{X}/D, \bar{\xi}_{1}) \longrightarrow
  \pi_{1}^{t}(X_{o}/D'_{o}, \bar{\xi}_{o}) \longrightarrow e ,
\]
\[
  \mu^{\infty}_{X_{o},\,\mathrm{mod}}(\bar{\xi}_{o}) \longrightarrow
  I_{o} \longrightarrow 0 ,
\]
où $\mu^{\infty}_{X_{o},\,\mathrm{mod}} = \varprojlim_{n} \mu_{n}$, la limite
étant prise sur les $n$ \emph{inversibles dans $\mathcal{O}_{X_{o}}$} --- donc,
si $p$ est la caractéristique de $X_{o}$, un groupe isomorphe à
$\prod_{\ell \neq p} \mathbb{Z}_{\ell}$ après choix d'une
orientation.\footnote{Le manuscrit écrit l'indice \emph{tame} en anglais ; il
est rendu ici \emph{mod}, pour \emph{modéré}, conformément à l'usage du reste
du carton. La restriction aux $n$ inversibles sur $X_{o}$ est portée en marge
de la page 87 et elle est essentielle : c'est elle qui fait de
$\mu^{\infty}$ un objet \emph{modéré}, et c'est elle qui, à l'exemple II
ci-dessous, l'annulera sur $\operatorname{Spec}\mathbb{Z}$.}

Le terme $\pi_{1}^{t}(X_{o}/D'_{o})$ s'identifie à $\pi_{1}^{t}(\mathcal{U})$
et à $\pi_{1}^{t}(\mathcal{X}/D')$ : c'est le théorème de spécialisation, écrit
à la volée sous la suite. Autrement dit : \emph{ôter la fibre spéciale ne change
pas le groupe fondamental modéré}, et tout le contenu du calcul est dans
l'extension par $I_{o}$, qui est la monodromie autour de cette fibre.

L'extension est \emph{centrale} et donc déterminée par une classe
\[
  \alpha \;\in\; H^{2}\bigl(\pi_{1}^{t}(\mathcal{U}),\, I_{o}\bigr) ,
\]
une fois choisie une orientation de $I_{o}$ via celles de
$\pi_{1}^{t}(\mathcal{U})$.

\subsection*{Lemme 1 : le revêtement universel modéré n'a pas de $H^{1}$}

\textbf{Lemme 1.} Soit $\tilde{\mathcal{U}}^{t}$ le revêtement universel de
$\mathcal{U}$ défini par $\pi_{1}^{t}$. Alors
$H^{1}(\tilde{\mathcal{U}}^{t}, \mathcal{J}) = 0$ pour $\mathcal{J}$ constant,
premier aux caractéristiques considérées.

C'est ce qui rend la suite spectrale de Hochschild--Serre utilisable en bas
degré : elle donne alors
\[
  0 \to H^{2}\bigl(\pi_{1}^{t}(\mathcal{U}), I\bigr)
  \xrightarrow{\ \rho\ } H^{2}(\mathcal{U}, I) \to
  H^{2}(\tilde{\mathcal{U}}^{t}, I)^{\pi_{1}^{t}} \to
  H^{3}\bigl(\pi_{1}^{t}, I\bigr) ,
\]
de sorte que $\alpha$ est connu dès qu'on connaît $\rho(\alpha)$ et qu'on sait
que son image plus loin est nulle. Et, toujours par le lemme 1 --- c'est
l'argument de Hurewicz ---,
\[
  H^{2}(\tilde{\mathcal{U}}^{t}, I)^{\pi_{1}^{t}} \;=\;
  \operatorname{Hom}\bigl(\pi_{2}(\tilde{\mathcal{U}}^{t}), I\bigr)^{\pi_{1}^{t}}
  \;=\; \operatorname{Hom}\bigl(\pi_{2}^{t}(\mathcal{U}), I\bigr)^{\pi_{1}^{t}} .
\]
C'est pourquoi $\pi_{2}$ apparaît ici : il mesure exactement ce qui empêche la
classe $\alpha$ d'être lue sur $H^{2}$ de $\mathcal{U}$.\footnote{Le
$\pi_{2}$ dont il s'agit est celui que la page 104 définit, dans le lot 6 :
pour $X$ simplement connexe, $M \rightsquigarrow H^{2}(X,M)$ est exact donc
pro-représentable, et $\pi_{2}(X) $ est l'objet qui le pro-représente ; dans le
cas général on prend le revêtement universel. C'est le $\pi_{2}$ de l'homotopie
étale, systématisé par Artin et Mazur en 1969 ; le carton l'utilise ici comme
un objet acquis.}

\subsection*{Lemme 2 : la classe est celle de l'auto-intersection}

\textbf{Lemme 2.} Soit $i : X_{o} \to \mathcal{X}$ l'inclusion et soit $\beta$
la classe de l'auto-intersection de $X_{o}$ dans $\mathcal{X}$, c'est-à-dire la
classe du fibré normal
\[
  \mathcal{L}_{X_{o}/\mathcal{X}} \;\simeq\; i^{*}\bigl(\mathcal{O}(X_{o})\bigr)
\]
dans $\mathrm{Pic}(X_{o}) = H^{1}(X_{o}, \mathbb{G}_{m})$. Soit
$\sigma : H^{1}(X_{o}, \mathbb{G}_{m}) \to H^{2}(\mathcal{U}, I)$
l'homomorphisme canonique --- image de Kummer suivie de restriction. Alors
\[
  \rho(\alpha) \;=\; \sigma(\beta) .
\]

\emph{C'est le théorème.} Il dit que la monodromie modérée autour de la fibre
spéciale est déterminée par le fibré normal de cette fibre --- c'est-à-dire par
son nombre d'auto-intersection, si $X_{o}$ est une courbe dans une
surface.\footnote{La page 89 écrit « $\beta \in H^{1}(X_{o}, \mu)$ », la page
95 « $\alpha \in H^{2}(X_{o}, \mu^{\infty}_{X_{o},\mathrm{tame}})$ » pour la
même quantité. Les deux écritures sont correctes et différentes : la classe
d'un fibré en droites est dans $H^{1}(\mathbb{G}_{m}) = \mathrm{Pic}$, sa
classe de Chern dans $H^{2}(\mu_{n})$, et la suite de Kummer
$1 \to \mu_{n} \to \mathbb{G}_{m} \to \mathbb{G}_{m} \to 1$ envoie la première
sur la seconde. La lecture ci-dessus fait la distinction que les deux pages ne
font pas.}

Grothendieck ajoute la réserve qu'il faut : cela détermine l'extension
\emph{sous réserve} de contrôler le passage de $I$ à $\mu$, c'est-à-dire de
savoir que $\mu_{o} \to I_{o}$ ne perd rien sur l'image de
$\pi_{2}^{t}(\mathcal{U})$. Il pose $I'_{o} = \mu_{o}/\operatorname{Im}
\pi_{2}^{t}(\mathcal{U})$ et observe que l'image de $\beta$ tombe alors dans
\[
  \ker\Bigl(H^{2}(\mathcal{U}, I'_{o}) \to
  H^{2}(\tilde{\mathcal{U}}^{t}, I'_{o})\Bigr) \;\simeq\;
  H^{2}\bigl(\pi_{1}^{t}, I'_{o}\bigr) ,
\]
donc définit bien une extension.

\pagerange{91}{93}
\section*{Le cas où la classe s'annule : Picard formel (pages 91 et 93)}

Les pages 91 et 93 sont écrites d'une plume très fine sur un papier fortement
traversé par le tapuscrit du verso, et une bonne part n'en est pas lisible.
L'articulation est néanmoins claire, et le calcul du milieu de la page 93 est
complet.

Il s'agit du cas où la classe $\beta$ est dans le noyau de
$H^{2}(\mathcal{X}, \mu) \to H^{2}(\mathcal{U}, \mathcal{J}_{o})$. Supposons
$\mathcal{J}_{o} \simeq \mu_{n}$ avec $n$ inversible sur $X_{o}$. L'hypothèse
dit alors que l'image de
$\beta_{o} \in \mathrm{Pic}(X_{o}) = H^{1}(X_{o}, \mathbb{G}_{m})$ dans
$H^{2}(X_{o}, \mu_{n})$ est nulle, c'est-à-dire que
\[
  \beta_{o} \;=\; n\gamma_{o} \;+\; \sum_{i} \lambda_{i}\,\beta_{o}^{i} ,
\]
où les $\beta^{i}$ sont les classes des autres composantes du diviseur.

Il faut relever cette relation de $X_{o}$ à $\mathcal{X}$, et c'est ici que le
schéma formel travaille. Soit $X_{m}$ le $m$-ième voisinage infinitésimal de
$X_{o}$. Alors :

\begin{quote}
le noyau de $\mathrm{Pic}(X_{m}) \to \mathrm{Pic}(X_{o})$ est \emph{uniquement
divisible par $n$}.
\end{quote}

La raison, que la page ne donne pas, est immédiate et vaut d'être écrite : ce
noyau est filtré par les $H^{1}(X_{o}, \mathcal{I}^{k}/\mathcal{I}^{k+1})$, qui
sont des groupes de cohomologie de $\mathcal{O}_{X_{o}}$-modules ; et la
multiplication par $n$ est bijective sur un $\mathcal{O}_{X_{o}}$-module
puisque $n$ est inversible dans $\mathcal{O}_{X_{o}}$.\footnote{C'est
exactement l'endroit où l'hypothèse de \emph{modération} paie : la restriction
« $n$ inversible sur $X_{o}$ » de la page 87, qui semblait n'être qu'une
définition, est ce qui rend ce noyau uniquement divisible et permet le
relèvement. Sans elle, tout le raisonnement s'arrête ici. La page l'assène
(« on voit que le noyau \dots\ est toujours uniquement divisible par $n$ »)
sans la relier à sa définition.}

Il en va donc de même du noyau de
\[
  \mathrm{Pic}(\mathcal{X}) \;=\; \varprojlim_{m} \mathrm{Pic}(X_{m})
  \longrightarrow \mathrm{Pic}(X_{o}) ,
\]
d'où la surjectivité de ${}_{n}\mathrm{Pic}(\mathcal{X})$ sur la $n$-torsion du
quotient, et l'existence d'un $\gamma \in \mathrm{Pic}(\mathcal{X})$ relevant
$\gamma_{o}$, avec
\[
  \beta \;=\; n\gamma + \sum_{i} a_{i}\beta_{i}
  \qquad \text{dans } \mathrm{Pic}(\mathcal{X}) .
\]

\pagerange{95}{97}
\section*{Le théorème, et deux exemples (pages 95 et 97)}

\subsection*{L'énoncé}

Une fois obtenue une relation $n\gamma = \sum_{i} c_{i}\beta_{i}$ sur
$\mathcal{X}$, on applique le \emph{théorème de Kummer} pour $\mathcal{X}/D$ :
il fournit un revêtement de $\mathcal{X}$, modéré le long de $D_{o}$, galoisien
de groupe $\mathcal{J}_{o} \simeq \mu_{n}$, trivial au-dessus de $X_{o}$. On
obtient ainsi la suite exacte complète
\[
  \pi_{2}^{t}(X_{o}/D'_{o}, \bar{\xi}_{o}) \to
  \mu^{\infty}_{X_{o},\,\mathrm{mod}}(\bar{\xi}_{o}) \to
  \pi_{1}^{t}(\mathcal{X}/D, \bar{\xi}_{1}) \to
  \pi_{1}^{t}(X_{o}/D'_{o}, \bar{\xi}_{o}) \to 0 ,
\]
et la conclusion :

\begin{quote}
$\pi_{1}^{t}(\mathcal{X}/D, \bar{\xi}_{1})$, comme extension de
$\pi_{1}^{t}(X_{o}/D'_{o}, \bar{\xi}_{o})$, est connu \emph{entièrement} en
termes de $X_{o}/D'_{o}$ et de la classe
$\alpha \in H^{2}\bigl(X_{o}, \mu^{\infty}_{X_{o},\,\mathrm{mod}}\bigr)$,
qui est l'auto-intersection de $X_{o}$ dans $\mathcal{X}$.
\end{quote}

\subsection*{Exemple I : une courbe sur un corps algébriquement clos}

Soit $X_{o}$ une courbe sur $k$ algébriquement clos. Alors
$\pi_{2}^{t}(X_{o} - D'_{o}) = 0$ \emph{sauf} si $X_{o} \simeq
\mathbb{P}^{1}_{k}$ et $D'_{o} = \varnothing$, auquel cas
$\pi_{2}^{t}(X_{o}) = \mu^{\infty}_{X_{o},\,\mathrm{mod}}$. Deux cas, donc, et
ils sont opposés.

\begin{enumerate}
  \item[1)] Si $X_{o} - D'_{o} \neq \mathbb{P}^{1}_{k}$ : alors
    $\pi_{1}^{t}(\mathcal{X}/D, \bar{\xi}_{1})$ est une extension
    \emph{triviale} de $\pi_{1}^{t}(X_{o} - D_{o}, \bar{\xi}_{o})$ par
    $\mu^{\infty}_{\mathrm{mod}}(\bar{\xi}_{o})$. La monodromie et le groupe
    fondamental de la fibre se séparent, et il n'y a rien à calculer.
  \item[2)] Si $X_{o} - D'_{o} = \mathbb{P}^{1}_{k}$, c'est-à-dire
    $X_{o} = \mathbb{P}^{1}_{k}$ et $D'_{o} = \varnothing$ : alors
    \[
      \pi_{1}^{t}(\mathcal{X}/D, \bar{\xi}_{1}) \;\simeq\; \mu_{n}(k) ,
    \]
    où $n$ est la partie première à $p = \operatorname{car} k$ du degré de
    l'auto-intersection $X_{o} \cdot X_{o}$ dans $\mathcal{X}$.
\end{enumerate}

Le second cas est celui qu'il faut retenir, et c'est le plus concret du carton.
Une droite projective d'auto-intersection $-d$ dans une surface fibrée a, pour
groupe fondamental modéré de son voisinage formel épointé, le groupe des
racines $n$-ièmes de l'unité, $n$ étant la partie de $d$ première à la
caractéristique. La monodromie est cyclique d'ordre $d$ : c'est ce que l'on
attend d'une singularité quotient cyclique, et le calcul le rend.

\subsection*{Exemple II : les entiers}

\begin{enumerate}
  \item[a)] $X_{o} = \operatorname{Spec}\mathbb{Z}$. Alors
    $\mu^{\infty}_{\mathrm{mod},\,X_{o}} = 0$, donc
    \[
      \pi_{1}^{t}(\mathcal{X}/D, \bar{\xi}_{1}) \;\simeq\;
      \pi_{1}^{t}(X_{o}/D'_{o}, \bar{\xi}_{o}) :
    \]
    il n'y a pas de monodromie du tout.\footnote{La page l'affirme sans dire
    pourquoi, et la raison est la définition même : $\mu^{\infty}_{\mathrm{mod}}$
    est la limite des $\mu_{n}$ pour $n$ \emph{inversible dans}
    $\mathcal{O}_{X_{o}}$ ; sur $\operatorname{Spec}\mathbb{Z}$, aucun entier
    $n > 1$ n'est inversible, et la limite est prise sur la seule valeur
    $n = 1$. C'est l'exemple qui montre le mieux ce que « modéré » interdit.}
    Si $D'_{o} = \varnothing$, on trouve $0$ --- c'est le théorème de
    Minkowski : $\mathbb{Q}$ n'a pas d'extension non ramifiée. Si
    $D'_{o} \neq \varnothing$, la partie abélienne se calcule par la théorie du
    corps de classes et vaut
    \[
      \prod_{\ell \in D'_{o}} (\mathbb{Z}/\ell\mathbb{Z})^{*} ,
    \]
    les extensions abéliennes de $\mathbb{Q}$ modérément ramifiées en les seuls
    $\ell \in D'_{o}$ étant les sous-corps de
    $\mathbb{Q}(\zeta_{\prod \ell})$. Il ajoute que le groupe entier doit être
    fini par dévissage, et pose une question qu'il laisse ouverte sur les
    corps de nombres modérément ramifiés sur $\mathbb{Q}$.
  \item[b)] $X_{o} = \operatorname{Spec}\mathbb{Z} - S$, avec
    $S \neq \varnothing$. Alors
    \[
      \mu^{\infty}_{\mathrm{mod},\,X_{o}} \;=\;
      \prod_{\ell \in S} \mu^{\infty}(\ell)_{X_{o}} ,
    \]
    puisque les $n$ inversibles sur $\mathbb{Z}[1/S]$ sont exactement ceux dont
    les facteurs premiers sont dans $S$. Et la page se termine sur une
    question, à laquelle rien ne répond : « Mais $\pi_{2}^{t}(X_{o}/D'_{o})$
    ??? ».
\end{enumerate}

\pagerange{99}{100}
\section*{Abhyankar semi-global, et la même classe une deuxième fois (pages 99 et 100)}

Deux feuillets non paginés ouvrent une question nouvelle, et le second
réénonce, sous forme de programme, le théorème qu'on vient de démontrer.

\subsection*{Une version semi-globale du résultat d'Abhyankar}

Soit $X$ un schéma localement noethérien connexe et $D = \sum D_{i}$ un
diviseur à croisements normaux à composantes régulières. Supposons
$Z = \bigcap_{i} D_{i} \neq \varnothing$ et le couple $(X, Z)$
\emph{hensélien}. Posons $\mathcal{U} = X - D$ et soit $\xi$ un point
géométrique de $\mathcal{U}$. Alors :

\begin{enumerate}
  \item[a)] $\mathcal{U}$ est connexe, et
    $\pi_{1}(\mathcal{U})_{\mathrm{mod}} \to \pi_{1}(X)$ est surjectif ;
  \item[b)] pour chaque $D_{i}$, le groupe d'inertie relatif à $D_{i}$ est un
    sous-groupe distingué \emph{commutatif} $I_{i}$ de
    $\pi_{1}(\mathcal{U})_{\mathrm{mod}}$, et l'on a un épimorphisme canonique
    \[
      \prod_{\ell} T_{\ell}(\Omega^{*}) \longrightarrow I_{i} ,
    \]
    le produit étant pris sur les $\ell$ distincts de la caractéristique en
    $D_{i}$ ;
  \item[c)] les $I_{i}$ commutent entre eux, et leur \emph{produit} est le
    noyau de
    $\pi_{1}(\mathcal{U}, \xi)_{\mathrm{mod}} \to \pi_{1}(X, \xi)$.
\end{enumerate}

L'énoncé a) est immédiat. Pour b), l'invariance des $D_{i}$ dans un revêtement
modérément ramifié $X'$ de $X$ est de nature étale locale et $X'$ est connexe
donc irréductible, ce qui identifie $I_{i}$ ; le fait qu'il soit quotient du
produit des modules de Tate est le calcul local classique. Pour c), la
commutation des $I_{i}$ est locale en un point du
croisement.\footnote{Cette description --- le groupe fondamental modéré du
complémentaire d'un diviseur à croisements normaux dans un couple hensélien
est une extension de $\pi_{1}(X)$ par un produit de modules de Tate, un par
composante --- est le calcul qui, en géométrie logarithmique, donne le groupe
fondamental \emph{log}-étale d'un point logarithmique. Le carton en fait ici
une « version semi-globale du résultat d'Abhyankar » ; c'est le même énoncé.}

\subsection*{Le programme kummérien : la classe, une deuxième fois}

La page 100 reprend, pour un $X$ quelconque, ce que le lemme 2 démontrait pour
un schéma formel.

Soient $X$ un schéma, $(D_{i})_{1 \leqslant i \leqslant r}$ des diviseurs de
Cartier $\geqslant 0$, $(n_{i})$ des entiers $\geqslant 1$,
$\mu_{\underline{n}} = \prod_{i} \mu_{n_{i}}$ sur $X$, $\Gamma$ un schéma en
groupes, et une extension de groupes sur $X$
\[
  1 \longrightarrow \mu_{\underline{n}} \longrightarrow G \longrightarrow
  \Gamma \longrightarrow 1 .
\]
Donnons-nous $X' \to \tilde{X}_{1} \to X$ avec $G$ opérant sur $\tilde{X}_{1}$,
$\mu_{\underline{n}}$ y opérant trivialement, de sorte que $\tilde{X}_{1}$
devienne un \emph{torseur} sous $\Gamma$ ; et $X'$ kummérien au-dessus de
$\tilde{X}_{1}$ relativement aux $D_{i}$, c'est-à-dire
\[
  X' \;=\; \operatorname{Spec}
  \mathcal{O}_{\tilde{X}}[T_{1}, \dots, T_{r}]\,\big/\,
  \bigl(T_{1}^{n_{1}} - a_{1},\ \dots,\ T_{r}^{n_{r}} - a_{r}\bigr) .
\]
L'extension $G$ de $\Gamma$ par $\mu_{\underline{n}}$ définit une classe
$\alpha \in H^{2}(B_{\Gamma}, \mu_{\underline{n}})$, et le torseur
$\tilde{X}$ fournit un homomorphisme
\[
  \varphi : H^{2}(B_{\Gamma}, \mu_{\underline{n}}) \longrightarrow
  H^{2}(X, \mu_{\underline{n}}) \;\simeq\;
  \prod_{i} H^{2}(X, \mu_{n_{i}}) .
\]

\begin{quote}
\textbf{Ce qu'on veut prouver.} Si $\varphi(\alpha) = (\xi_{i})_{1 \leqslant i
\leqslant r}$, alors
\[
  \xi_{i} \;=\; c\bigl(\mathcal{O}_{X}(D_{i})\bigr) \pmod{n_{i}} .
\]
\end{quote}

C'est le lemme 2 de la page 89, écrit sans schéma formel : \emph{la classe de
l'extension d'inertie est la classe de Chern du diviseur}. La page 102, dans le
lot 6, en donne une rédaction en corollaire ; et la page 200, tout à la fin du
carton, l'écrit une quatrième fois, pour la gerbe des racines
$\underline{n}$-ièmes des $D_{i}$ : « L'existence d'une section dépend de la
nullité d'un élément de $H^{2}(X, \mu_{\underline{n}})$, qui n'est autre,
comme on devine, que la classe de cohomologie des $D_{i}$ ». Quatre énoncés,
cent dix pages, une seule identité.

La page se clôt sur une ligne qui indique où cela devait aller : « Variantes
\dots\ formelle, ou \dots\ loc.\ annelés ».

\end{document}
